\subsection{Grafy}
Pro vizualizaci aktivit vývojářů byla zvolena knihovna Highcharts, jež byla do prostředí frameworku Vue.js integrována prostřednictvím rozšiřujícího modulu \texttt{highcharts-vue}. V rámci implementace byly vytvořeny dvě specializované komponenty: \texttt{LineChart} a \texttt{PolarChart}. Tyto komponenty zajišťují zapouzdření logiky vykreslování a oddělují prezentační vrstvu od datové.

Vzhled a chování jednotlivých grafů je definováno prostřednictvím konfiguračních objektů. Samotná data jsou komponentám předávána dynamicky z příslušného controlleru. Pro zajištění konzistence zobrazení a minimalizaci duplicity kódu dochází uvnitř komponent ke sloučení těchto vstupních dat s výchozím nastavením. Ukázka implementace komponenty pro liniový graf, demonstrující dynamické skládání konfiguračního objektu, je uvedena na výpisu kódu \ref{listing:highcharts-example}.

Dalším z důležitých aspektů implementace byla podpora světlého a tmavého režimu, která je realizována pomocí modulu \texttt{highcharts/themes/adaptive}. Tento modul zajišťuje automatické přizpůsobení barevného schématu grafů aktuálnímu nastavení aplikace. Detekce preferovaného režimu probíhá na základě systémové preference (pomocí CSS pravidla \texttt{prefers-color-scheme}), případně dle CSS třídy, která je dynamicky přidávána na HTML element \texttt{body} na~základě uživatelského nastavení.

\begin{listing}[h!]
    \caption{Vue.js komponenta s konfigurací grafu}\label{listing:highcharts-example}
    \begin{minted}{vue}
<script setup lang="ts">
import { merge } from 'lodash';
import { computed } from 'vue';

const props = defineProps<{
    options: object;
}>();

const baseOptions = {
    chart: { height: 250 },
    legend: {
        align: 'right',
        verticalAlign: 'top',
    },
    // ...
    title: { text: null },
    xAxis: {
        dateTimeLabelFormats: {
            day: '%d. %m.',
            week: '%d. %m.',
        },
        lineColor: '#CCD6EB',
        tickColor: '#CCD6EB',
        type: 'datetime',
    },
    yAxis: {
        minTickInterval: 1,
        title: { enabled: false },
    },
};

const finalOptions = computed(() => {
    return merge({}, baseOptions, props.options);
});
</script>

<template>
    <highcharts :options="finalOptions"></highcharts>
</template>
    \end{minted}
\end{listing}
